\documentclass[11pt,a4paper]{article}
\usepackage[top=2.5cm,bottom=2.5cm,left=2.5cm,right=2.5cm]{geometry}
\usepackage[T1]{fontenc}
\usepackage[utf8]{inputenc}
\usepackage{lmodern}
\usepackage{amsmath}
\usepackage{enumitem}
\usepackage{listings}
\usepackage{xcolor}
\usepackage[most]{tcolorbox}

% ── Code style ────────────────────────────────────────────────────────────────
\lstdefinestyle{code}{
  basicstyle=\ttfamily\small,
  backgroundcolor=\color{gray!10},
  frame=single,
  breaklines=true,
  tabsize=4,
  showstringspaces=false,
}

% ── Coloured boxes ────────────────────────────────────────────────────────────
\tcbset{
  basebox/.style={
    enhanced, breakable, sharp corners,
    boxrule=0pt, leftrule=3pt,
    left=10pt, right=10pt, top=8pt, bottom=8pt,
    fonttitle=\bfseries\small,
    attach title to upper,
    after title={\par\smallskip\hrule\smallskip},
  },
  setupbox/.style={basebox, colback=gray!8,  colframe=black!60, coltitle=black!80},
  infobox/.style ={basebox, colback=blue!5,  colframe=blue!50!black, coltitle=blue!50!black},
  notebox/.style ={basebox, colback=yellow!8, colframe=orange!70!black, coltitle=orange!70!black},
}

\setlength{\parskip}{6pt}
\setlength{\parindent}{0pt}

% ──────────────────────────────────────────────────────────────────────────────
\begin{document}

\begin{center}
  {\Large\bfseries ME403 --- Lab 1: Warming up to the Robot!}\\[4pt]
  {\normalsize Sabanci University, Spring 2025-26 \quad|\quad Due: April 8--9}
\end{center}

\rule{\linewidth}{0.8pt}

Please complete the following tasks on a robot of your choice (Dobot Magician or MG400).

% ── Assignment ────────────────────────────────────────────────────────────────
\section*{Assignment}

\textbf{Task 1.} Derive the formulation of the end-effector position in the robot's frame by
using homogeneous transformation matrices of the links. Then, please write a function that
utilizes these matrices to estimate the tip position. All the required dimensions such as link
lengths and relative axis locations can be obtained by simply measuring them at one of the
robots, or by referring to the robot manual, which can be found online. This part of your code
should estimate where the robot's tip would end up in the Cartesian space when it is at a
certain configuration (i.e.\ for certain joint angles).

\textbf{Task 2.} Instruct the robot to move to 20 random different configurations within its
workspace. At each configuration collect the position feedback provided by the robot and store
that data in an array. Compare those feedback values against your code's estimation. If your
forward kinematics model is alright, your estimation and the values returned from the robot
should have a very good match. If you see significant errors, please review your kinematic
model. Show the agreement between measured and estimated position values by plotting the
magnitude of the difference between them.

Let $e = p - p_m$, where $p = [x\ y\ z]^T$ is the estimated position and
$p_m = [x_m\ y_m\ z_m]^T$ is the position provided by the robot. Please present the magnitude
of the error in those 20 configuration samples with a box plot.

\textbf{Task 3.} Write a function that obtains the Jacobian of the end-effector and outputs the
matrix to the terminal window. This function should get the current configuration as input and
provide the matrix at will.

\rule{\linewidth}{0.4pt}

% ── Due date ──────────────────────────────────────────────────────────────────
\section*{Due Date \& Deliverables}

\textbf{Due Date:} Lab sessions on April 8--9.

\textbf{Deliverables:} Please complete these steps on your own by the time of your lab session.
A brief report will be requested after the lab session. An assignment for that post-lab document
will be announced later.

\rule{\linewidth}{0.4pt}

% ── Provided files ────────────────────────────────────────────────────────────
\section*{Provided Files}

\begin{tabular}{ll}
  \texttt{utils.py}         & Robot connection and motion helpers. \\
  \texttt{interface.py}     & Terminal interface to jog the robot and run your code. \\
  \texttt{myCode.py}        & \textbf{Write your code here} (inside \texttt{run()}). \\
  \texttt{requirements.txt} & Python dependencies. \\
\end{tabular}

% ── Quick start ───────────────────────────────────────────────────────────────
\section*{Quick Start --- If You Attended Lab 0}

Your virtual environment is already set up. Re-activate it and install the requirements.
 \textbf{Do this at the start of every lab}, as we may update the packages between sessions.

\begin{tcolorbox}[setupbox, title={Linux / macOS}]
\begin{lstlisting}[style=code]
# From Students/01_SecondWeek/
source .venv/bin/activate
pip install -r requirements.txt
\end{lstlisting}
\end{tcolorbox}

\begin{tcolorbox}[setupbox, title={Windows (PowerShell)}]
\begin{lstlisting}[style=code]
# From Students\01_SecondWeek\
.\.venv\Scripts\Activate.ps1
pip install -r requirements.txt
\end{lstlisting}
\end{tcolorbox}

% ── Robot selection ───────────────────────────────────────────────────────────
\section*{Select Your Robot}

Open \texttt{utils.py} and set the robot type at the top before running anything:

\begin{lstlisting}[style=code]
ROBOT_TYPE = "magician"     # or "mg400"
MG400_IP   = "192.168.2.x"  # change to your assigned robot's IP (MG400 only)
\end{lstlisting}

\begin{tcolorbox}[infobox, title={If you are using the Dobot Magician}]
Check the \textbf{``Working with Dobot Magician''} file on Sucourse before working
with the Magician.
\end{tcolorbox}

% ── How to run ────────────────────────────────────────────────────────────────
\section*{How to Run}

\begin{lstlisting}[style=code]
python interface.py
\end{lstlisting}

At the prompt:
\begin{itemize}[nosep]
  \item \texttt{move q1 q2 q3 q4} --- jog the robot to the given joint angles
  \item \texttt{x} --- execute your \texttt{run()} function in \texttt{myCode.py}
  \item \texttt{q} --- quit and close the connection
\end{itemize}

% ── First-time setup (missed Lab 0) ───────────────────────────────────────────
\newpage
\rule{\linewidth}{1pt}
\begin{center}
  {\large\bfseries First-Time Setup --- If You Did Not Attend Lab 0}
\end{center}
\rule{\linewidth}{1pt}

\begin{tcolorbox}[notebox, title={Note: Avoid PDF Copy-Paste for Commands}]
Type commands manually or copy from the digital version. Some PDF viewers insert
extra spaces in folder paths (e.g.\ \texttt{.\ venv} instead of \texttt{.venv}).
\end{tcolorbox}

\subsection*{1 --- Virtual Environment}

\begin{tcolorbox}[setupbox, title={Linux / macOS}]
\begin{lstlisting}[style=code]
# From Students/01_SecondWeek/
python3 -m venv .venv
source .venv/bin/activate
pip install -r requirements.txt
\end{lstlisting}
\end{tcolorbox}

\begin{tcolorbox}[setupbox, title={Windows (PowerShell)}]
\begin{lstlisting}[style=code]
# From Students\01_SecondWeek\
python -m venv .venv
.\.venv\Scripts\Activate.ps1
pip install -r requirements.txt
\end{lstlisting}
\end{tcolorbox}

\subsection*{2 --- Dobot Magician (USB/Serial)}

The Magician connects over USB via a CP210x Silicon Labs chip at 115\,200 baud.
\texttt{pydobotplus} handles the serial protocol; the firmware runs on the robot's
internal controller. Close DobotStudio / DobotDemo before running any Python script
--- only one process can hold the serial port at a time.

\textit{Linux only} --- grant port access once per machine, then log out and back in:
\begin{lstlisting}[style=code]
sudo usermod -a -G dialout $USER
ls /dev/ttyUSB*    # verify: should show /dev/ttyUSB0
\end{lstlisting}

\textit{Windows} --- the CP210x driver installs automatically in most cases. If not,
download it from Silicon Labs. The port shows up as \texttt{COM3} (or similar) in
Device Manager under \textit{Ports (COM \& LPT)}.

\subsection*{3 --- DOBOT MG400 (Ethernet/TCP-IP)}

The MG400 does not use serial; it communicates over Ethernet via three TCP ports:

\begin{tabular}{lll}
  Port 29999 & Dashboard & Enable/disable, error query, pose and angle readout \\
  Port 30003 & Move API  & MovJ, MovL, Arc, RelMovJ, Sync, \ldots \\
  Port 30004 & Feedback  & 1440-byte packets every 8\,ms (live robot state) \\
\end{tabular}

\vspace{6pt}

\textit{Network} --- set your PC's Ethernet adapter to static IP \texttt{192.168.2.100},
netmask \texttt{255.255.255.0}, then ping your assigned robot:
\begin{lstlisting}[style=code]
ping 192.168.2.7    # Robot 1 -- replace with your assigned robot's IP
\end{lstlisting}

\textit{Robot IP addresses:}

\begin{tabular}{ll}
  Robot 1 & \texttt{192.168.2.7}  \\
  Robot 2 & \texttt{192.168.2.10} \\
  Robot 3 & \texttt{192.168.2.9}  \\
  Robot 4 & \texttt{192.168.2.6}  \\
\end{tabular}

\vspace{6pt}

\textit{SDK} --- the MG400 requires Dobot's Python SDK, which is not on PyPI.
Clone it once from the \texttt{dobot\_ws/} root:
\begin{lstlisting}[style=code]
git clone https://github.com/Dobot-Arm/TCP-IP-4Axis-Python.git \
    vendor/TCP-IP-4Axis-Python
\end{lstlisting}

\texttt{utils\_mg400.py} picks up \texttt{dobot\_api.py} from \texttt{vendor/}
automatically --- no further configuration needed.

\end{document}
